%% =====================================================================
%%  10 Appendix
%% =====================================================================
\appendix

\section{The no-bend identity}
\label{app:nobend}

Write the velocity of the ball in components, $\vect v=(v_x,v_y,v_z)$, and
let the aerodynamic force be parallel to $\vect v$ (drag only). Then the
horizontal part of the equation of motion reads
\begin{equation}
  \dot v_i = -k_d\,|\vect v|\,v_i, \qquad i=x,y ,
  \label{eq:app-drag}
\end{equation}
so that
\begin{equation}
  \frac{\dd}{\dd t}\!\left(\frac{v_y}{v_x}\right)
  = \frac{\dot v_y v_x - \dot v_x v_y}{v_x^{\,2}}
  = \frac{-k_d|\vect v|\,(v_y v_x - v_x v_y)}{v_x^{\,2}} = 0 .
  \label{eq:app-identity}
\end{equation}
The ratio $v_y/v_x$ is therefore a constant of the motion, and so is the
ratio of the horizontal displacements,
\begin{equation}
  \frac{y(t)-y_0}{x(t)-x_0} = \frac{v_{y0}}{v_{x0}} ,
  \label{eq:app-ratio}
\end{equation}
for any drag law whose direction is $-\vect v$, whatever the speed
dependence $k_d(|\vect v|)$ is. Setting $x=x_g=\SI{11}{m}$,
$y_0=0$, $x_0=0$ gives
\begin{equation}
  y_{\rm entry} = x_g\,\frac{v_{y0}}{v_{x0}} ,
  \label{eq:app-nobend}
\end{equation}
which does not contain $C_D$, the density, the mass, or the spin decay. This
is Eq.~\eqref{eq:nobend} of the main text. A force with a component
perpendicular to $\vect v$ breaks the identity: the Magnus force contributes
$+k_m\,\vect\omega\times\vect v$ to the acceleration, and the right-hand side
of \eqref{eq:app-identity} becomes
$k_m[(\vect\omega\times\vect v)_y v_x-(\vect\omega\times\vect v)_x v_y]/v_x^{2}$,
which is not zero because $\vect\omega\times\vect v$ is perpendicular to
$\vect v$. The measured
entry ordinate $y=\SI{2.50}{m}$ therefore differs from
\eqref{eq:app-nobend} only because of lift; the whole lateral answer is a
lift measurement.

\section{Uncertainty propagation for the kinematic estimator}
\label{app:jacobian}

With $\vect\theta=(v_{x0},v_{y0},v_{z0},a_x,a_y,a_z)$ and observations
$\vect r_1=\vect r(t_1)$, $\vect r_2=\vect r(t_2)$, the model is linear in
$\vect\theta$, so the least-squares estimate is exact. Writing
$\Delta\vect r_k=\vect r_k-\vect r_0$,
\begin{equation}
  \vect r_k = t_k\,\vect v_0 + \tfrac12 t_k^{2}\,\vect a ,
  \qquad k=1,2 ,
\end{equation}
the Jacobian with respect to the data is block-structured,
\begin{equation}
  \mathbf{J} =
  \begin{pmatrix}
    t_1\mathbf{I} & \tfrac12 t_1^{2}\mathbf{I}\\[2pt]
    t_2\mathbf{I} & \tfrac12 t_2^{2}\mathbf{I}
  \end{pmatrix}
  \in\mathbb R^{6\times 6} ,
\end{equation}
and for independent per-coordinate noise the covariance of the
estimate is $\mathbf{C}=\sigma_{\rm pos}^{2}(\mathbf{J}^{\!\top}\mathbf{J})^{-1}$.
For the prediction $g=\bigl(t_{\rm entry},y_{\rm entry},z_{\rm entry}\bigr)$
the gradient $\mathbf{G}=\partial g/\partial\vect\theta$ is evaluated by central
differences on the closed-form quadratic of Section~\ref{sec:method}; the
reported standard deviations are
$\sigma_g=\sqrt{\operatorname{diag}(\mathbf{G}\mathbf{C}\mathbf{G}^{\!\top})}$.
The largest entries of $\mathbf{C}$ are
$\sigma_{a_y}=\sigma_{a_z}=\sqrt{6}\,\sigma_{\rm pos}/\Delta t^{2}$, which
for $\sigma_{\rm pos}=\SI{10}{mm}$ and $\Delta t=\SI{0.05}{s}$ gives
$\SI{9.8}{m/s^2}$ --- the number quoted in Section~\ref{sec:results} and the
reason the acceleration is only marginally determined.

\section{Priors, likelihood and sampling}
\label{app:priors}

The main line fits six parameters,
$\vect\theta=(v_{x0},v_{y0},v_{z0},C_D,\beta_1,\beta_2)$, where
$\vect\beta=\beta_1\vect e_1+\beta_2\vect e_2$ is expanded in the two unit
vectors perpendicular to the launch velocity, Eq.~\eqref{eq:beta}. The priors
are independent normal distributions,
\begin{align}
  \vect v_0 &\sim \mathcal N\!\left((30,0,7),\,
     \operatorname{diag}(15^{2},15^{2},10^{2})\right)\ \si{m/s}, \\
  C_D &\sim \mathcal N(0.25,\,0.15^{2}), &
  \beta_i &\sim \mathcal N(0,\,0.008^{2})\ \si{m^{-1}},
     \quad i=1,2 ,
\end{align}
truncated by $5<v_{0x}<60$, $|v_{0y}|<30$, $0<v_{0z}<30$, $0<C_D<1$ and
$|\beta_i|<0.08\ \si{m^{-1}}$, as listed in Table~\ref{tab:priors}. The
likelihood is Gaussian in the six position residuals at
$\sigma_{\rm pos}=\SI{5}{mm}$ per coordinate,
\begin{equation}
  -2\ln\mathcal L = \sum_{k=1}^{2}
  \frac{\lVert \vect r_k-\vect r(t_k;\vect\theta)\rVert^{2}}{\sigma_{\rm pos}^{2}},
\end{equation}
with $\vect r(t;\vect\theta)$ obtained by integrating
Eq.~\eqref{eq:eom} with the classical fourth-order Runge--Kutta scheme at
$\Delta t=\SI{5}{ms}$ up to the observation epochs and at
$\Delta t=\SI{2.5}{ms}$ for the crossing. The maximum a posteriori point is
found by a damped Gauss--Newton iteration on the joint residual vector of
Eq.~\eqref{eq:residual}: six data residuals, each divided by $\sigma_{\rm pos}$,
together with the six prior pulls. The Laplace covariance
$\mathbf{C}=(\mathbf J^{\mathsf T}\mathbf J)^{-1}$, inflated by a factor
$1.15$ and evaluated at the map, defines a Student-$t$ proposal with eight
degrees of freedom; every retained draw is re-weighted by the exact
prior-times-likelihood ratio, Eq.~\eqref{eq:weights}, so the reported
distribution is not the proposal but the posterior. The production run keeps
$131\,072$ weighted draws with an effective sample size of $109\,585$; it is
the run whose quantiles appear in Tables~\ref{tab:intervals}
and~\ref{tab:scenarios}. Section~\ref{sec:posterior} quotes the same run.

\section{Reproduction script}
\label{app:code}

The listing below reproduces every number quoted in
Sections~\ref{sec:results} and~\ref{sec:conclusion} from the two given
positions. It reads the two numerical modules shipped with the paper
(\texttt{bkick.py}, the integrator, and \texttt{estimators.py}, the two
estimators); it is included from the file that produced the numbers so the
printed listing cannot drift from the code that was run.

\begin{lstlisting}[language=Python,caption={Core of the reproduction script
\texttt{code/predict.py} (abridged to the two estimator calls and the outcome
count; the full file ships with the paper and prints every number quoted in
Sections~\ref{sec:results} and~\ref{sec:conclusion}).}]
import numpy as np, bkick as B, estimators as E
R_OBS = np.array([[1.50, 0.45, 0.44], [2.96, 0.86, 0.74]])   # m
T_OBS, SIGMA = np.array([0.05, 0.10]), 0.005                  # s, m
e1 = E.e1_predict(R_OBS, T_OBS, SIGMA, n_mc=40000, seed=0)     # constant a
post = E.PhysicalPosterior(R_OBS, T_OBS, SIGMA)               # six parameters
th, _ = post.map()
out = post.sample(n=40000, burn=8000, thin=2, seed=1)
t, y, z = post.predict(out["chain"])
print(E.summarize(t, y, z))          # mean, sd, quantiles, P(clean goal)
\end{lstlisting}

\section{Marginal intervals and the joint region}
\label{app:joint}

The paper quotes two marginal $90\%$ intervals, one for $y$ and one for $z$, and
these are not a $90\%$ joint region. For an approximately Gaussian predictive
density the distinction is elementary but worth making explicit, because the two
constructions differ by a factor that grows with the correlation and with the
dimension. Let $\vect\Sigma$ be the $2\times2$ predictive covariance and
$\chi^2_{2}(p)$ the chi-square quantile with two degrees of freedom. The
rectangle enclosed by the two marginal intervals contains
\begin{equation}
  P_{\rm rect}=\Pr\!\left(|y-\bar y|\le z_{1-p/2}\sigma_y,\;
                        |z-\bar z|\le z_{1-p/2}\sigma_z\right),
  \label{eq:rect}
\end{equation}
which for independent coordinates with $p=0.10$ is $0.90^2=0.81$ and rises
towards $0.90$ only as the correlation goes to unity: two marginal $90\%$
intervals contain between $81$ and $90\%$ of the probability, and two marginal
$95\%$ intervals between $90$ and $95\%$. The ellipse
$\vect\Delta^{\top}\vect\Sigma^{-1}\vect\Delta\le\chi^2_2(0.90)=4.605$ is the
region that actually carries $90\%$. We therefore report the interval of each
coordinate separately and never add the two numbers up as if they bounded a
region; where a joint statement is needed, as for the clearance probability of
Eq.~\eqref{eq:clearance}, the joint density is used directly.

The crossing time is itself uncertain, and propagating it linearly would mix two
error sources. The trajectory is advanced on a fixed step, so the crossing is
located by a quadratic interpolation of $x(t)$ across the bracketing step. With
step $h$ and offset $T\in[0,h)$ the interpolant is
\begin{equation}
  x(t)\simeq x_k+\dot x_kT+\tfrac12\ddot x_kT^2,\qquad T=t-t_k,
\end{equation}
and the two Lagrange weights that propagate a position error into the crossing
coordinates are $L_1=T(2h-T)/h^2$ and $L_2=T(T-h)/2h^2$, so that the entry-point
error is $\vect e=L_1\vect r_1+L_2\vect r_2$ and its variance is
$(L_1^2+L_2^2)\sigma_{\rm pos}^2$. Because $L_1+L_2\ne1$ the two fixes do not
average; the closer one carries almost all of the weight, which is why a crossing
near the last fix is far more sensitive to that fix. Our analytic Jacobian of the
crossing coordinates agrees with centred finite differences to
$1.6\times10^{-8}$, and a $10^6$-sample perturbation of the two fixes reproduces
the quoted $\sigma_y$ and $\sigma_z$ to better than $1\%$.

\section{The whole-sphere clearance criterion}
\label{app:sphere}

A shot passes cleanly if the whole ball, not its centre, crosses the plane inside
the frame. Writing the aperture as the rectangle $|y|\le b$, $0\le z\le h$ with
$b=\SI{3.66}{m}$, $h=\SI{2.44}{m}$ and radius $R=\SI{0.110}{m}$, the ball clears
if the centre lies in the inset rectangle $|y|\le b-R$, $R\le z\le h-R$, which is
what Eq.~\eqref{eq:clearance} evaluates. The inset test is exact for a rectangular
aperture and a spherical ball, because the distance from the centre to each edge
is measured perpendicular to that edge. The two cases where it needs care are
the posts: a ball whose centre passes outside $|y|=b$ but whose surface still
touches the post is neither clean nor wide, and we count it as clipping, as the
outcome table of \S\ref{sec:answer} states. When a trajectory is sampled rather
than integrated to the plane, the same test must be applied to the segment
between two samples: with sampling step $h$ and speed $|v|$ the surface can cross
an edge between samples, so the clearance test is applied on the segment,
\begin{equation}
  d(t)^2=\big|\vect\delta_0+\vect u\,t\big|^2,\qquad
  \vect\delta_0=\vect c_0-\vect c_{\rm edge},\quad \vect u=\vect v\,h,
\end{equation}
a quadratic in $t$ whose minimum on $[0,1]$ is checked analytically. The
closed-form shortcut $g/\sqrt{1+r^2}-R$, with $r$ the slope of the trajectory in
the $y$--$z$ plane and $g$ the distance from the centre to the nearest edge, is
valid only when the perpendicular foot of the trajectory onto that edge lies
inside the edge segment; otherwise the nearest point is a corner and the
shortcut overestimates the clearance. All probabilities quoted in this paper use
the exact segment test.
